% HEADLINE: Applications Engineer -- ARM SoCs, NPUs, industrial Ethernet, camera sensors
% Persona D (FAE / Solutions Engineer), 2 pages. Every number cites a Bullet Bank ID in a trailing % comment.
% \textbf{⚠️ ID} = unreconciled value (Open Reconciliations), must be replaced before a final build — never pick a value.
% Arguments: (1) evaluates and designs in vendor technology, (2) builds reference pipelines and tools other teams adopt, (3) works with customers, factories and partners.
% Customer-facing framing: each Wenglor bullet names the vendor stack and what was handed to the next team (PoC, process, tool, guidance).
% Ortiq shown (persona D rule) as a short entry; sorted by end date (Wenglor = Present > Ortiq Oct 2026 > Yandex). RoboCup is an Experience/leadership entry here.
% No W-01/W-03/W-04 used. Kazan title as on the cv; RaccoonLab = "Principal Engineer" [stated 2026-10-29].

\section{Summary}
\cvitem{}{Embedded engineer who has sat on the customer side of five vendor stacks -- Hilscher netX, WIBU CodeMeter, AMLogic NPU, Sony IMX, AMD Zynq -- evaluated them and built the PoCs. Brings hands-on Linux depth, tools that other teams adopt, and factory-floor work with 6 OEMs.}   % W-12 W-13 W-02 W-07 (5 named stacks) · W-08 (2 adopting teams) · Y-01 (6 OEMs; ⚠️ R11 is on the project count only, not used here) — 44 words

\section{Experience}

\cventry{Nov 2023 -- Present}{Embedded Linux Software Developer}{Wenglor MEL GmbH}{Unterschleißheim, Germany}{}{   % F-01 ⚠️ R4: cv title kept until R4 is ticked
\begin{itemize}
  \item Prototyped EtherNet/IP on \textbf{Hilscher CiFX / netX90} and evaluated PROFINET and EtherCAT; defined the integration path for the sensor lines.   % W-12 (Arg 1) — 3 protocols
  \item Led the software-licensing infrastructure end to end on \textbf{WIBU CodeMeter} (architecture, integration, rollout); Wenglor now earns revenue by selling sensor feature licenses.   % W-13 (Arg 1) — [stated 2026-10-29]; no revenue/SKU/headcount numbers (⚠️ R33 R34 R35)
  \item Integrated the AMLogic A311D NPU via OpenCL (30~ms to 3~ms per frame) and Sony IMX296/IMX567 sensors on A311D and Zynq UltraScale+.   % W-02 (✅ R3) + W-07 (2 sensors, 2 SoCs) — Arg 1
  \item Designed and shipped a C++ test and production tool for the B60 smart-camera line, concept to deployment, adopted by 2 teams.   % W-08 (Arg 2) — production + development teams
  \item Architected a modular Ethernet vision server (nginx- and GenICam-inspired) now reused across 3 product families.   % W-06 (Arg 2) — 3 families
  \item Wrote the Cyber Resilience Act path: read-only rootfs and A/B updates shipped, container isolation for user apps proposed.   % W-11 (Arg 2) — 2 shipped + 1 proposal
  \item Ran a development-process workshop for the firmware team; mentored engineers and students on V4L2, CRA and HIL testing.   % W-17 (Arg 3) — 1 workshop
\end{itemize}}

\cventry{Jun 2026 -- Oct 2026}{Founder \& CTO}{Ortiq}{Munich, Germany}{}{Built a decision-intelligence / spatial-analytics SaaS from zero: engine, web app, infra, compliance. \textbf{⚠️ O-01}}   % O-01 (Arg 3) — ⚠️ R15: no reconciled metric; replace with one metric or delete the placeholder. "3 startups" (O-02, ⚠️ R16) NOT used. Description form to save a line.

\cventry{Jan 2023 -- Sep 2023}{OEM Process Manager}{Yandex, Smart Devices -- Factory Team}{Moscow, Russia}{}{
\begin{itemize}
  \item Synchronised production plans for \textbf{⚠️ Y-01} smart-device projects across 6 Shenzhen OEMs, each line at ~7,000~units/day per colour.   % Y-01 (Arg 3) — 6 OEMs, 7,000/day ✅; project count ⚠️ R11 (9 vs 9+)
  \item Built one Python toolchain (OpenHTF, ADB, U-Boot, Ansible/Salt) automating 4 factory workflows: flashing, SN/MAC, HDCP keys, QA.   % Y-02 (Arg 3) — 4 workflows
  \item Re-shaped QA/QC/AQL with the OEM partners and introduced metrics-driven defect tracking.   % Y-04 (Arg 3) — qualitative; scope = OEM partners
\end{itemize}}

\cventry{Feb 2021 -- Jan 2023}{Embedded Development Lead}{Kazan Scientific Center, Quantum Optics Lab}{Kazan, Russia}{}{
\begin{itemize}
  \item Designed the Spatial Light Modulator control protocol and public SDK (Cortex-M firmware, Linux drivers over USB and MIPI CSI-2), used across several optical experiments.   % K-01 (Arg 2) — scope: several experiments
  \item Published the SDK with Sphinx docs and containerised CI (GitHub Actions, Docker) so other groups could build it.   % K-02 (Arg 2) — qualitative
\end{itemize}}

\cventry{Aug 2019 -- Jan 2021}{Principal Engineer}{RaccoonLab (Innopolis)}{Innopolis, Russia}{}{   % [stated 2026-10-29] org + title; dates ⚠️ R24
\begin{itemize}
  \item Developed autonomous-drone CAN modules end to end -- Altium hardware, FreeRTOS firmware, DroneCAN/Cyphal protocol, bench tooling -- for PX4/ArduPilot airframes.   % I-06 + I-03 (Arg 1, lead) — no count: ⚠️ R19 open, non-numeric form
  \item Took the modules into production: ran the ramp-up, costed each batch and qualified substitutes for scarce components.   % I-07 (Arg 3, design-in angle; absorbs I-05's substitution half) — no units/customers: ⚠️ R20 R22 open
  \item Co-invented a registered visual-positioning system for autonomous multirotor landing (patent RU 2022610258).   % I-02 — 1 patent. I-01 ANPR (96.75 %) now only in Papers (⚠️ R32: RaccoonLab or university work?)
\end{itemize}}

\cventry{2014}{Team Captain, PCB \& Firmware}{R2T robotics team}{Mashhad, Iran}{}{
\begin{itemize}
  \item Led the team to the RoboCup 2014 world title (Junior Rescue B, João Pessoa) and the IranOpen 2014 national title.   % E-01 (Arg 3, earliest leadership) — 1st place ×2; ⚠️ R13 "2nd then 1st" NOT used; "Team Captain" user-stated 2026-10-29
\end{itemize}}

\cventry{2013 -- 2019}{Embedded \& PCB Developer}{Freelance and FUM Robotics Lab}{Mashhad, Iran}{}{
\begin{itemize}
  \item Shipped firmware and PCBs for 4 client products: dental X-ray controller, kiln thermometer, A/C servicing machine, STM32 LoRaWAN/UWB nodes.   % E-02 — 4 shipped boards
\end{itemize}}
